% ECONOMICS_PROBLEM: {"id": "L15-269", "title": "Portability and interoperability", "jel_code": "L15", "source_id": "OP-0409"}
% !TeX program = xelatex
\documentclass[11pt,a4paper]{article}
\usepackage[margin=23mm]{geometry}
\usepackage{fontspec}
\setmainfont{Latin Modern Roman}
\usepackage{amsmath,amssymb,mathtools}
\usepackage{xcolor,enumitem,booktabs,longtable,array,xurl,hyperref}
\definecolor{muted}{HTML}{53636C}
\hypersetup{colorlinks=true,urlcolor=blue,linkcolor=blue}
\setlength{\parindent}{0pt}
\setlength{\parskip}{5pt}
\newcommand{\R}{\mathbb R}
\newcommand{\E}{\mathbb E}
\newcommand{\N}{\mathbb N}
\newcommand{\ind}{\mathbf 1}
\newcommand{\Var}{\operatorname{Var}}
\newcommand{\Cov}{\operatorname{Cov}}
\newcommand{\supp}{\operatorname{supp}}
\DeclareMathOperator*{\argmin}{arg\,min}
\DeclareMathOperator*{\argmax}{arg\,max}
\newcommand{\entrymeta}[3]{{\sffamily\small\color{muted}\textbf{\texttt{#1}}\quad|\quad #2\par #3\par}}
\newcommand{\Problemref}[1]{\texttt{#1}}
\newcommand{\JELref}[2]{\texttt{#2} (JEL \texttt{#1})}
\begin{document}
\section*{Portability and interoperability}
\textbf{Problem ID:} \texttt{L15-269}\quad \textbf{JEL:} \texttt{L15}\par
% BEGIN ECONOMICS STATEMENT
\label{problem:OP-0409}
{\sffamily\small\color{muted}Original subject group: Digital, platform and data economics\par}
\label{definitions:problem:30007691}
\textbf{Audit classification:} Background; proposed extension. Metadata corrected to distinguish first-online 2025 from volume publication 2026. The cited paper already proves special-case welfare results, and the joint investment frontier is editorial.
\subsection*{Definitions and precise research target}
Consider two competing social platforms with heterogeneous users, switching costs and network benefits. Let policy \(r=(p,i,s)\) specify portability \(p\), interoperable messaging \(i\), and minimum security standard \(s\), each chosen from a declared finite menu. Platforms invest in useful data production and security, price access, and choose quality after the policy is announced. Users then choose a platform or multihome. For a specified distribution \(\theta\) of preferences, switching costs, technology costs and breach losses, let \(\mathcal E(r,\theta)\) be the equilibrium set under a stated timing protocol. Define \(W(e)\) as consumer surplus plus platform profits, with profits measured net of real investment expenditure and surplus net of expected breach harm in equilibrium \(e\). Let \(I(e)\) be data-production investment and \(B(e)\) expected breach harm. The design task is to compare feasible policies by expected \(W(e)\) under an explicit equilibrium-selection rule, subject to selection-weighted expected investment \(I_r\geq I_{\min}\) and expected breach loss \(B_r\leq B_{\max}\), where the two bounds are externally specified policy requirements. Determine robustness to alternative switching-cost and breach-risk estimates. This proposed extension does not claim that interoperability is always beneficial; existing stylized welfare results are acknowledged, while the joint investment-and-security design frontier lacks verified open-statement provenance.

\textbf{Definitions, assumptions and mathematical qualifications.} Portability is transfer of specified user records; interoperability permits specified cross-platform communication; neither implies the other. Give distinct policy menus, implementation costs and security-investment constraints. A complete two-platform game specifies user types, multi-homing rules, demand, prices, productive-data investments, breach probabilities, losses and timing. Restrict policies to those with nonempty equilibrium sets and fix a supported selection kernel. Let \(I_r=E_{\kappa_r}I(e)\), \(B_r=E_{\kappa_r}B(e)\), \(W_r=E_{\kappa_r}W(e)\); expected constraints are \(I_r\geq I_{\min},B_r\leq B_{\max}\). Requiring them for every selected equilibrium is a stronger variant. Check that a policy satisfies both bounds before asking for its welfare maximizer.
\subsection*{Variants and assumption changes}
\begin{enumerate}
\item \textbf{Sourced benchmark.} Kim analyzes a specific two-period Hotelling model and breach liability. Its welfare conclusions hold under that model assumptions; they do not establish universal interoperability benefits.
\item \textbf{Editorial extension.} Permit endogenous productive-data investment, heterogeneous users and separate portability/interoperability choices; identify or estimate the constrained welfare frontier, keeping breach-risk limits explicit.
\end{enumerate}
\subsection*{References and source audit}
\textbf{Checked source.} Journal of Economics \& Management Strategy 35(2) (2026), pp. 219--232, Sections 2--6; first online 21 May 2025. This article proves specific portability/interoperability welfare and liability results; it does not establish the proposed frontier as open. Full primary text retrieved. \url{https://onlinelibrary.wiley.com/doi/full/10.1111/jems.12643}.
\begin{enumerate}\raggedright
\item\hypertarget{definitions:source:30007691:1}{} Jeong-Yoo Kim. 2026. \emph{Data Portability and Interoperability Between Digital Platforms}. Journal of Economics \textbackslash{}\& Management Strategy 35(2) (2026), pp. 219--232, Sections 2--6; first online 21 May 2025.
\url{https://onlinelibrary.wiley.com/doi/full/10.1111/jems.12643}.
DOI: \url{https://doi.org/10.1111/jems.12643}.
\end{enumerate}

\subsection*{Common conventions: Source mathematical conventions}

These conventions apply to each entry unless explicitly replaced.

\textbf{Models and identification.} A primitive model \(m\) specifies all probability laws, feasible choices, information, equilibrium and measurement rules used by its target. The admissible class \(\mathcal M\) is fixed independently of outcomes. If \(P_O(m)\) is its observable probability law and \(\tau(m)\) the target, its identified set at observable law \(P\) is
\[
 \mathcal I_\tau(P)=\{\tau(m):m\in\mathcal M,\ P_O(m)=P\}.
\]
Point identification means that this set is a singleton. An empty set signals incompatibility of data and assumptions. Coordinate bounds need not describe the joint set. Bounds are sharp only if every retained target is attainable or arbitrarily approachable by compatible models. Finite bounds require explicit restrictions on unobserved quantities; unrestricted confounding, hidden wealth, extrapolation or arbitrary equilibrium selection can yield uninformative sets.

\textbf{Causal interventions.} A potential outcome \(Y_i(p)\) is the outcome under a complete policy regime \(p\), including timing, financing, information and market spillovers. The target population and units are fixed before treatment. With interference, use the complete assignment vector or a justified exposure mapping; individual no-interference notation alone cannot identify an economy-wide intervention. A difference of mean potential outcomes does not require the cross-world joint distribution. Conditioning on post-treatment migration, survival, enrollment or employment changes the estimand and needs separate justification.

\textbf{Statistical statements.} An estimator, confidence set, forecast or test specifies a sampling law, model class and loss/coverage criterion. Uniform \((1-\alpha)\)-coverage means \(\inf_{m\in\mathcal M}\Pr_m\{\tau(m)\in C_n\}\geq1-\alpha\), or an explicitly asymptotic version. The whole parameter space trivially has coverage, so informativeness requires a stated width, risk or power objective. A forecasting improvement is relative to a named benchmark, information set and evaluation population.

\textbf{Equilibrium and optimization.} An equilibrium comprises feasible plans, optimality, beliefs where needed, budget constraints and all market/resource clearing. The primitives or a stated benchmark define each correspondence \(\mathcal E(p;m)\). Nonemptiness is assumed or proved, never inferred just from compact policy sets. A selected-equilibrium target uses a declared selection rule; a robust target quantifies over all permitted equilibria. Use suprema or infima unless attainment is established. An \(\varepsilon\)-optimal policy is feasible and within \(\varepsilon\) of the finite optimum value. Report an empty feasible set separately.

\textbf{Welfare and accounting.} Normative utility, interpersonal normalization, social weights, numeraire, discounting and the use of public receipts are primitives. Behavioral utility need not coincide with the welfare criterion. Household welfare includes ownership income and rents once; adding profits again to compensating variation generally double-counts them. A monetary equivalent is defined by an explicit transfer and price convention, with monotonicity and range conditions. Average effects, mechanism contributions, transfers and welfare losses are different targets. Infinite sums require convergence and dynamic feasible plans require terminal or no-Ponzi conditions.

\textbf{Source versions.} A working-paper date, revision date and journal publication year are distinguished. A DOI for a journal version is not automatically the DOI of an earlier manuscript. Locators refer to the stated version and printed pages where specified. A metadata-only check does not verify a claim about a full-text conclusion. Every entry's source subsection narrows the provenance claim accordingly.
% END ECONOMICS STATEMENT
\end{document}
